\documentclass[11pt]{article}

\usepackage[a4paper,margin=1in]{geometry}
\usepackage{amsmath,amssymb,amsthm,mathtools}
\usepackage{booktabs}
\usepackage{array}
\usepackage{enumitem}
\usepackage{xcolor}
\usepackage{hyperref}

\hypersetup{
    colorlinks=true,
    linkcolor=blue,
    urlcolor=blue,
    citecolor=blue
}

\newtheorem{theorem}{Theorem}[section]
\newtheorem{lemma}[theorem]{Lemma}
\newtheorem{corollary}[theorem]{Corollary}
\newtheorem{proposition}[theorem]{Proposition}
\theoremstyle{remark}
\newtheorem{remark}[theorem]{Remark}

\DeclareMathOperator*{\argmax}{arg\,max}
\DeclareMathOperator{\ext}{ext}

\title{\textbf{The Vertex-Localization Threshold Is Zero}}
\author{Formal Results V.1--V.14}
\date{\today}

\begin{document}

\maketitle

\begin{abstract}
The default composite objective has a vertex maximizer at zero interaction
strength, but that maximizer is first-order degenerate along the
$(\lambda,\sigma)$ face. The \texttt{face\_bowl} interaction introduces a
positive inward derivative along both of these directions. Consequently, every
positive interaction strength, however small, moves the maximizer into the
interior of the face. The true vertex-localization threshold is therefore
\[
s^*=0.
\]

The displacement from the vertex is linear in $s$, while the objective-value
gap is quadratic:
\[
\delta(s)=\frac{3}{8}s+O(s^2),
\qquad
\Delta(s)=\frac{9}{32}s^2+O(s^3).
\]
Coarse grid searches can miss this failure because the displacement and value
gap become small faster than practical grid resolution can detect them.
\end{abstract}

\tableofcontents

\section{Executive Summary}

The objective considered in this report is defined on the box
\[
H=[0,1]\times[0,1]\times[0,2]\times[0,3],
\]
with coordinates $(\lambda,\sigma,b,k)$.

The separable variables $b$ and $k$ are always maximized at their upper
bounds. The remaining optimization is over the $(\lambda,\sigma)$ face.

At $s=0$, the vertex
\[
\theta_c=(1,0,2,3)
\]
is a maximizer. However, the inward derivatives in the $\lambda$ and
$\sigma$ directions are both zero. The vertex is therefore not
first-order robust.

The \texttt{face\_bowl} perturbation gives each of these directions an inward
slope of $3s/4$. Thus, for every $s>0$, an inward move increases the objective.
The vertex immediately ceases to be a local maximizer.

The main conclusions are
\[
\boxed{s^*=0},
\]
\[
\boxed{\delta(s)=\frac{3}{8}s+O(s^2)},
\]
and
\[
\boxed{\Delta(s)=\frac{9}{32}s^2+O(s^3)}.
\]

The broader methodological conclusion is that a vertex-only optimization
strategy is reliable only after checking both:

\begin{enumerate}[label=\arabic*.]
    \item that the base objective is globally maximized at the candidate vertex;
    \item that the vertex has sufficient inward margin and curvature under the
    intended perturbations.
\end{enumerate}

\section{Setup}

Consider the objective
\begin{align}
C_s(\theta)
={}& b+k+
\left[1-(1-\lambda)^2\right]
+\left[1-\sigma^2\right]
\nonumber\\
&+
s\left(1-\left(\lambda-\frac12\right)^2\right)
\left(1-\left(\sigma-\frac12\right)^2\right),
\end{align}
where $s\geq 0$ and
\[
\theta=(\lambda,\sigma,b,k)\in H.
\]

The variables $b$ and $k$ occur only in the separable term $b+k$. Therefore,
\[
b^*=2,\qquad k^*=3
\]
at every maximizer.

The problem consequently reduces to the $(\lambda,\sigma)$ face.

Introduce centered coordinates
\[
u=\lambda-\frac12,
\qquad
w=\frac12-\sigma.
\]
Then
\[
u,w\in\left[-\frac12,\frac12\right].
\]

The sign reversal in the definition of $w$ maps the relevant corner
\[
(\lambda,\sigma)=(1,0)
\]
to
\[
(u,w)=\left(\frac12,\frac12\right).
\]

Using
\[
1-(1-\lambda)^2
=
\frac34+u-u^2
\]
and
\[
1-\sigma^2
=
\frac34+w-w^2,
\]
the reduced objective becomes
\begin{align}
C_s
&=
\frac{13}{2}
+
u+w-u^2-w^2
+
s(1-u^2)(1-w^2).
\end{align}

The additive constant $13/2$ does not affect the location of the maximizer. Define
\[
f_s(u,w)
=
u+w-u^2-w^2+s(1-u^2)(1-w^2).
\]

\section{Reduction to a Symmetric Face}

\begin{lemma}[Reduction to the symmetric face]
The variables $b$ and $k$ attain their upper bounds at every maximizer:
\[
b^*=2,\qquad k^*=3.
\]
The remaining objective is
\[
f_s(u,w)
=
u+w-u^2-w^2+s(1-u^2)(1-w^2),
\]
on
\[
\left[-\frac12,\frac12\right]^2.
\]
Moreover, $f_s$ is symmetric under the exchange $u\leftrightarrow w$.
\end{lemma}

\begin{proof}
The first claim follows because $b+k$ is strictly increasing in both $b$ and $k$. The symmetry follows directly from
\[
f_s(u,w)=f_s(w,u).
\]
\end{proof}

Once uniqueness of the maximizer has been established, this symmetry will imply that the maximizer lies on the diagonal $u=w$.

\section{Strict Concavity}

\begin{theorem}[Strict concavity]
For every $s\geq 0$, the reduced objective $f_s$ is strictly concave on
\[
\left[-\frac12,\frac12\right]^2.
\]
\end{theorem}

\begin{proof}
The Hessian is
\[
\nabla^2 f_s(u,w)
=
\begin{pmatrix}
-2-2s(1-w^2) & 4suw\\
4suw & -2-2s(1-u^2)
\end{pmatrix}.
\]

Therefore,
\[
-\nabla^2 f_s(u,w)
=
\begin{pmatrix}
2+2s(1-w^2) & -4suw\\
-4suw & 2+2s(1-u^2)
\end{pmatrix}.
\]

The diagonal entries are at least $2$. Since
\[
|u|,|w|\leq \frac12,
\]
we have
\[
|4suw|\leq s
\]
and
\[
1-u^2,\;1-w^2\geq \frac34.
\]

Consequently,
\begin{align}
\det(-\nabla^2 f_s)
&\geq
\left(2+\frac{3s}{2}\right)^2-s^2\\
&=
4+6s+\frac{5s^2}{4}>0.
\end{align}

Thus $-\nabla^2 f_s$ is positive definite, so $f_s$ is strictly concave.
\end{proof}

\begin{corollary}
The face objective has a unique global maximizer.
\end{corollary}

\begin{remark}
Strict concavity rules out multiple local optima. Therefore, an apparent
vertex/interior transition produced by a numerical grid must be distinguished
from the actual optimizer obtained analytically.
\end{remark}

\section{Exact Maximizer}

\begin{theorem}[Exact maximizer]
The unique maximizer is
\[
u^*=w^*=t^*(s),
\]
where $t^*(s)$ is the unique root in $[0,1/2]$ of
\[
2s\,t^3-2(1+s)t+1=0.
\]

At $s=0$,
\[
t^*(0)=\frac12.
\]
For every $s>0$,
\[
0<t^*(s)<\frac12.
\]
\end{theorem}

\begin{proof}
The stationarity equations are
\[
1-2u-2su(1-w^2)=0
\]
and
\[
1-2w-2sw(1-u^2)=0.
\]

Subtracting the two equations gives
\[
(u-w)\bigl[2+2s(1+uw)\bigr]=0.
\]

The bracket is strictly positive on the domain, so every stationary point
satisfies
\[
u=w=t.
\]

Substituting $u=w=t$ into either stationarity equation gives
\[
1-2t-2st(1-t^2)=0,
\]
or equivalently
\[
2s\,t^3-2(1+s)t+1=0.
\]

Strict concavity guarantees that the stationary point is the unique global
maximizer.
\end{proof}

In the original coordinates,
\[
\lambda^*=\frac12+t^*(s),
\qquad
\sigma^*=\frac12-t^*(s),
\qquad
b^*=2,
\qquad
k^*=3.
\]

Hence, for every $s>0$,
\[
\lambda^*<1,
\qquad
\sigma^*>0.
\]

The optimizer immediately leaves the vertex $(1,0,2,3)$.

\section{The Threshold Is Zero}

\begin{theorem}[Zero vertex-localization threshold]
The vertex
\[
\theta_c=(1,0,2,3)
\]
is optimal at $s=0$, but it ceases to be locally optimal for every $s>0$. Therefore,
\[
\boxed{s^*=0}.
\]
\end{theorem}

\begin{proof}
At the face corner $(u,w)=(1/2,1/2)$,
\[
\frac{\partial f_s}{\partial u}
=
1-2u-2su(1-w^2)
=
-\frac{3s}{4}.
\]

The inward direction is decreasing $u$. Therefore the inward directional
derivative is
\[
D_u^{\mathrm{in}}f_s
=
-\frac{\partial f_s}{\partial u}
=
\frac{3s}{4}.
\]

Similarly,
\[
D_w^{\mathrm{in}}f_s=\frac{3s}{4}.
\]

Both inward derivatives are strictly positive whenever $s>0$. Therefore an
inward move increases the objective, and the corner cannot be a local
maximizer for any positive $s$.
\end{proof}

At $s=0$, the inward derivatives in the $\lambda$ and $\sigma$ directions
vanish. By contrast, the $b$ and $k$ directions have strictly negative inward
derivatives, so those coordinates remain pinned at their upper bounds.

\section{Displacement and Objective Gap}

Define the displacement
\[
\delta(s)=\frac12-t^*(s).
\]

\begin{theorem}[Displacement and gap]
As $s\to0^+$,
\[
\delta(s)=\frac38s+O(s^2).
\]

Furthermore, if
\[
\Delta(s)
=
C_s(\theta^*)
-
\max_{\theta\in\ext(H)}C_s(\theta),
\]
then
\[
\Delta(s)
=
\frac{9}{32}s^2+O(s^3).
\]
\end{theorem}

\begin{proof}
Set
\[
t=\frac12-d.
\]

Substituting into the stationarity equation gives
\[
-\frac{3s}{4}
+
\left(2+\frac{s}{2}\right)d
+
3sd^2
-
2sd^3
=0.
\]

Since $d=O(s)$,
\[
d=\frac38s+O(s^2).
\]

Along the inward diagonal path
\[
u=w=\frac12-x,
\]
the exact objective difference from the corner is
\begin{align}
&f_s\left(\frac12-x,\frac12-x\right)
-
f_s\left(\frac12,\frac12\right)
\nonumber\\
&\qquad=
\frac{3s}{2}x
-\left(2+\frac{s}{2}\right)x^2
+sx^3-\frac{s}{2}x^4.
\end{align}

Maximizing this expression gives
\[
\Delta(s)=\frac{9}{32}s^2+O(s^3).
\]
\end{proof}

The exact negative Hessian at the corner is
\[
A=
-\nabla^2 f_s\left(\frac12,\frac12\right)
=
\begin{pmatrix}
2+\frac32s & -s\\
-s & 2+\frac32s
\end{pmatrix}.
\]

Its eigenvalue in the inward diagonal direction $(1,1)$ is
\[
2+\frac12s.
\]

Using the inward gradient
\[
\gamma=\frac{3s}{4}(1,1),
\]
the corresponding quadratic approximation to the gap is
\[
\Delta_{\mathrm{quad}}(s)
=
\frac{9s^2}{8(4+s)}.
\]

This is an approximation and should not be confused with the exact gap.

A useful rational approximation for the displacement is
\[
\delta_{\mathrm{Pad\acute e}}(s)=\frac{3s}{8+2s}.
\]

It matches the first-order expansion but is not exact for finite $s$.

\section{Why Grid Searches Missed the Failure}

Let $x$ denote the inward distance from the corner along the diagonal. The
exact crossing distance $x_c(s)$ is the positive root of
\[
\frac{3s}{2}
-\left(2+\frac{s}{2}\right)x
+sx^2-\frac{s}{2}x^3
=0.
\]

An interior grid point beats the corner whenever
\[
\frac{1}{n-1}<x_c(s).
\]

For small $s$,
\[
x_c(s)=\frac34s+O(s^2),
\]
so the required resolution behaves as
\[
n\approx 1+\frac{4}{3s}.
\]

This is an asymptotic estimate, not an exact identity involving
$2\delta(s)$.

The numerical difficulty has two sources:

\begin{enumerate}
    \item the optimizer displacement is only $O(s)$;
    \item the objective-value gap is only $O(s^2)$.
\end{enumerate}

Moreover, a four-dimensional grid with $n$ points per axis requires
\[
n^4
\]
objective evaluations.

For example, at $s=0.05$,
\[
\Delta(0.05)\approx 0.000694.
\]

A coarse grid can therefore report the corner as the apparent maximizer even
though the exact maximizer is interior.

A classification rule such as
\[
\texttt{localization\_at\_vertex}=(\texttt{gap}<0.05)
\]
is a tolerance-based numerical flag, not a mathematical certificate.

\section{A Corrected Local Robustness Criterion}

Let $x_i\geq0$ denote inward slack coordinates from a candidate vertex
$\theta_c$. Suppose the base objective has the local expansion
\[
C_0(\theta_c+x)
=
C_0(\theta_c)
-\sum_i a_i x_i
-\frac12x^\top Qx
+o(\|x\|^2),
\]
with $a_i\geq0$.

Suppose also that the perturbation satisfies
\[
P(\theta_c+x)
=
P(\theta_c)+\sum_i p_i x_i+\cdots.
\]

The perturbed first-order coefficient in coordinate $i$ is
\[
-a_i+sp_i.
\]

If $p_i>0$, the first-order stability limit on that axis is
\[
s_i=\frac{a_i}{p_i}.
\]

If $a_i=0$ and $p_i>0$, then
\[
s_i=0.
\]

This criterion is local. Before applying it, one must verify that $\theta_c$ is
actually a global maximizer of $C_0$.

\begin{remark}
A strict vertex maximizer need not have a positive first-order margin. For
example,
\[
C_0(x)=-x^2,\qquad x\geq0,
\]
has a strict maximum at $x=0$, although its inward derivative is zero.
Therefore, strict maximality and positive first-order robustness are distinct
properties.
\end{remark}

\section{Smooth Perturbations on a Quadratically Flat Face}

Suppose the base objective has local quadratic loss
\[
C_0(\theta_c+x)
=
C_0(\theta_c)-\frac12x^\top Qx+o(\|x\|^2),
\]
where $Q$ is positive definite.

Let the perturbation have a positive inward linear term
\[
P(\theta_c+x)
=
P(\theta_c)+g^\top x+o(\|x\|).
\]

The leading-order optimization problem is
\[
\max_{x\geq0}
\left\{
s\,g^\top x-\frac12x^\top Qx
\right\}.
\]

Ignoring active-cone constraints, the optimizer is
\[
x^*=sQ^{-1}g+O(s^2),
\]
and the gap is
\[
\Delta(s)
=
\frac12s^2g^\top Q^{-1}g+O(s^3).
\]

If $Q$ is diagonal and the perturbation is separable, this becomes
\[
\Delta(s)
=
s^2\sum_i\frac{\gamma_i^2}{2c_i}
+O(s^3).
\]

For coupled perturbations or coupled base curvature, the full quadratic form
must be retained.

\section{Fractional Perturbations}

Let the base loss be quadratic,
\[
-Ax^2,
\]
and suppose the perturbation behaves locally as
\[
P(x)-P(0)\sim\gamma x^\alpha,
\qquad 0<\alpha<2.
\]

The local objective is
\[
-Ax^2+s\gamma x^\alpha.
\]

Its optimizer satisfies
\[
x^*
=
\left(\frac{\alpha\gamma s}{2A}\right)^{1/(2-\alpha)},
\]
and the gap scales as
\[
\Delta(s)
=
\Theta\left(s^{2/(2-\alpha)}\right).
\]

The result applies only while the predicted optimizer remains in the region
where the local power-law approximation is valid.

\section{Coupled Perturbations and Directional Optimization}

For a coupled perturbation, coordinate-wise contributions can overestimate the
true gap. Let $d$ be an inward direction and write
\[
x=Rd.
\]

If the base loss is quadratic and $P$ is positively homogeneous of degree one,
then
\[
C_0(\theta_c+Rd)+sP(\theta_c+Rd)
=
C_0(\theta_c)
-\frac{R^2}{2}d^\top Qd
+sR\,D_dP+\cdots.
\]

Optimizing over $R$ gives
\[
\Delta(s)
=
s^2
\max_{d\in\mathcal K}
\frac{(D_dP)^2}{2d^\top Qd}
+o(s^2),
\]
where $\mathcal K$ is the inward direction cone.

This formulation allows for multiple maximizing directions. It does not require
the leading-order optimizer to lie on a unique ray.

For isotropic curvature $Q=cI$,
\[
\Delta(s)
=
\frac{s^2}{2c}
\max_{\substack{\|d\|=1\\d\in\mathcal K}}
(D_dP)^2
+o(s^2).
\]

\section{Amplitude Ceiling and Saturation}

Local derivative laws have a finite validity range. If the perturbation is
bounded, then
\[
\Delta(s)
\leq
s\left(\sup_H P-P(\theta_c)\right).
\]

If $P(\theta_c)=0$ and $P\geq0$, this simplifies to
\[
\Delta(s)\leq s\sup_H P.
\]

Any derivative-based prediction that exceeds this bound is invalid. The
perturbation has saturated before the local asymptotic optimizer is reached.

This issue is relevant for bounded ridge-like or angular interactions. A large
finite-difference slope near a corner does not imply an arbitrarily large
objective gain.

\section{General Flatness Law}

Suppose the base loss behaves as
\[
C_0(\theta_c)-C_0(\theta_c+x)
\sim A x^\beta,
\]
and the perturbation behaves as
\[
P(\theta_c+x)-P(\theta_c)
\sim \gamma x^\alpha,
\]
with
\[
0<\alpha<\beta.
\]

Then
\[
\max_x\{-Ax^\beta+s\gamma x^\alpha\}
\]
has optimizer
\[
x^*
=
\Theta\left(s^{1/(\beta-\alpha)}\right)
\]
and gap
\[
\boxed{
\Delta(s)
=
\Theta\left(s^{\beta/(\beta-\alpha)}\right).
}
\]

The quadratic-base result is obtained by setting $\beta=2$:
\[
\Delta(s)
=
\Theta\left(s^{2/(2-\alpha)}\right).
\]

A flatter base produces a smaller scaling exponent and therefore a larger gap at
small interaction strength.

\section{Weighted Anisotropic Law}

Suppose different coordinates have different base flatness orders:
\[
C_0(\theta_c)-C_0(\theta_c+x)
\asymp
\sum_i A_i x_i^{\beta_i}.
\]

Use the base-adapted dilation
\[
D_t x
=
\left(t^{1/\beta_i}x_i\right)_i.
\]

Suppose the leading perturbation is weighted-homogeneous:
\[
P(D_t x)-P(0)
=
t^q(P(x)-P(0))+o(t^q),
\qquad 0<q<1.
\]

Then
\[
\boxed{
\Delta(s)
=
\Theta\left(s^{1/(1-q)}\right).
}
\]

For a monomial
\[
P(x)=\gamma\prod_i x_i^{\alpha_i},
\]
the weighted degree is
\[
q=\sum_i\frac{\alpha_i}{\beta_i}.
\]

Thus anisotropic base flatness can change the scaling exponent even when the
ordinary degree of the perturbation remains unchanged.

\section{Base Self-Failure}

All preceding perturbation results assume that the candidate corner is the
global maximizer of the base objective $C_0$.

That assumption can fail independently of perturbations. For example,
\[
r(\lambda,\sigma)
=
-\bigl((1-\lambda)-0.25\bigr)^2
-(\sigma-0.35)^2
\]
has its maximum at
\[
(\lambda,\sigma)=(0.75,0.35),
\]
which is strictly interior.

A vertex search therefore returns the wrong point for every $s\geq0$. This is
not an $s^*=0$ perturbation phenomenon; it is a failure of the base model.

The appropriate diagnostic is global:
\[
\texttt{base\_self\_fails}
=
\left[
\max_H C_0
>
\max_{\theta\in\ext(H)}C_0
\right].
\]

This test should be performed before any local margin analysis.

\section{Implementation Implications}

The following changes should be reflected in the implementation and
documentation.

\subsection{Formal theorem documentation}

The vertex theorem may remain valid under its stated assumptions, but its
conclusion is not automatically robust. Under the default base objective, the
$(\lambda,\sigma)$ directions have zero first-order margin.

\subsection{\texttt{demonstrate\_nonlinear}}

The claim that ``small coupling preserves vertex localization'' should be
replaced by:

\begin{quote}
Small coupling can produce a small displacement and a small objective gap, but
it does not preserve the exact vertex maximizer when the base vertex is
first-order degenerate.
\end{quote}

\subsection{\texttt{localization\_at\_vertex}}

This should be documented as a resolution-and-tolerance flag rather than a
mathematical certificate.

The analysis should report:

\begin{itemize}
    \item the candidate vertex;
    \item whether the base objective self-fails;
    \item all inward first-order margins;
    \item local curvature or flatness order;
    \item predicted displacement scaling;
    \item predicted objective-gap scaling;
    \item numerical grid resolution;
    \item the perturbation amplitude ceiling.
\end{itemize}

\subsection{Recommended regression tests}

The test suite should verify:

\begin{enumerate}
    \item the sign of the inward directional derivative;
    \item the negative off-diagonal signs in $-\nabla^2f_s$;
    \item interior displacement for every $s>0$;
    \item
    \[
    \frac{\delta(s)}{s}\longrightarrow\frac38;
    \]
    \item
    \[
    \frac{\Delta(s)}{s^2}\longrightarrow\frac{9}{32};
    \]
    \item the distinction between the exact grid-crossing distance and
    $2\delta(s)$;
    \item base self-failure before perturbation analysis;
    \item coupled directional optimization;
    \item saturation against the amplitude bound.
\end{enumerate}

\section{Final Conclusions}

For the stated \texttt{face\_bowl} objective,
\[
\boxed{s^*=0}.
\]

The vertex is the optimizer only at $s=0$. Every positive interaction strength
moves the optimizer into the interior of the $(\lambda,\sigma)$ face.

The small-$s$ behavior is
\[
\boxed{
\delta(s)=\frac38s+O(s^2)
}
\]
and
\[
\boxed{
\Delta(s)=\frac{9}{32}s^2+O(s^3).
}
\]

The central result is mathematically sound after correcting the derivative
sign convention and the Hessian display. The report should also distinguish
carefully among:

\begin{itemize}
    \item exact global results;
    \item local asymptotic laws;
    \item heuristic approximants;
    \item numerical grid classifications.
\end{itemize}

The principal methodological lesson is:

\begin{quote}
Vertex localization is trustworthy only after global base validation and
local robustness analysis. A zero inward first-order margin means that an
arbitrarily small positive perturbation can move the optimizer away from the
vertex.
\end{quote}

\end{document}
